
\section{Case Study: What the Flags Actually Look Like}
\label{sec:flagcases}
Numbers alone undersell what the census found, so we open the flags
themselves. The 155 published flags are exact archive paths
(e.g., \texttt{NORMAL/IM-0133-0001.jpeg}); Table~\ref{tab:flagprops} compares their measured
acquisition properties against size-matched random controls
(Mann--Whitney U, exact values from
\texttt{flagged\_image\_properties.json}).
\begin{table}[h]
\centering\small
\begin{tabular}{lcccc}
\hline
Property & Flagged mean & Control mean & $U$ & $p$ \\
\hline
Sharpness (Laplacian var.) & 0.005301 & 0.003516 & 61109 & 1.74e-27 \\
Contrast (std of pixels) & 0.235458 & 0.218673 & 49102 & 4.93e-07 \\
Entropy (bits) & 7.320236 & 7.260964 & 44755 & 3.53e-03 \\
\hline
\end{tabular}
\caption{Flagged films are measurably blurrier, flatter and less
informative than controls. The sharpness separation
($p=1.7\times10^{-27}$) is the single most decisive statistic of the
forensics battery: label uncertainty concentrates on
acquisition-degraded films.}
\label{tab:flagprops}
\end{table}
Reading the first published flags --- \texttt{NORMAL/IM-0621-0001.jpeg}, \texttt{NORMAL/NORMAL2-IM-0587-0001-0002.jpeg} --- the
pattern a radiologist would call ``technically limited study'' is exactly
what the probe disagrees with: a NORMAL label on a film whose low
sharpness and contrast made the underlying report equivocal. The census
does not claim these films are certainly pneumonia; it claims their
labels carry substantially lower confidence than the dataset assumes,
which is the statement the +5.9-point cleaning gain then monetizes.

\section{Why the Gate is Necessary: A Failure Mode and Its Operational Fix}
\label{sec:gatetheory}
Confident learning estimates the joint distribution $Q_{\tilde y, y^*}$
of observed and latent labels from out-of-fold predicted probabilities.
Its correctness condition is that the out-of-fold probabilities are
rank-informative --- informally, that per-class probability mass exceeds
the self-confidence thresholds used for calibration. When the out-of-fold
model is at or below the majority-class rate, the estimated thresholds
collapse onto the class priors and the estimator returns a ``noise rate''
that is an artefact of the probe, not the data. We observed exactly this
fabrication on real collapse cases during development (a malaria probe
trained with too few steps: OOF accuracy below the majority rate, census
rate uninterpretable; the \texttt{tissuemnist} external probe likewise,
refused by the gate). The admissibility gate
\begin{equation}
\mathrm{OOF\ acc} \;\ge\; \max\!\big(0.60,\ \mathrm{majority} + 0.05\big)
\quad\wedge\quad \mathrm{steps}_{\mathrm{fold}} \ge 250
\end{equation}
is a sufficient operational surrogate for the rank-informativeness
condition. On this dataset the admitted probe ran at 89.81\%% OOF
accuracy against a 74.20\%% majority baseline --- a 15.6-point margin
over the gate's floor --- and its per-class thresholds (0.774, 0.845) sit
clear of the priors. The gate's cost is conservatism: genuinely noisy
but weakly-learnable datasets are refused rather than miscounted. For a
discovery claim --- ``an automated, falsifiable census of a benchmark''
--- conservatism
is the correct failure direction, and the gate is enforced inside the
shipped \texttt{dx-census} command rather than in the paper text alone.
